% Chapter 4, Section 4 _Linear Algebra_ Jim Hefferon
%  http://joshua.smcvt.edu/linearalgebra
%  2001-Jun-12
\section{Jordan 形式}
\index{Jordan form|(}
\noindent\textit{本节要用到三个选读小节中的内容：~子空间的组合、行列式的存在性
以及 Laplace 展开。}

本章开头我们回顾了每个线性映射 $\map{h}{V}{W}$ 都可以
关于某些基 $B\subset V$ 与 $D\subset W$
用一个部分单位矩阵来表示。
换一种说法，部分单位矩阵形式是矩阵等价的
一种规范形式。
本章考虑陪域等于定义域的情形，于是我们自然会问，
当两个基相等时，即有 $\rep{t}{B,B}$ 时，
能做到什么。
简言之，我们想要矩阵相似的规范形式。

我们已经指出，在 $B,B$ 的情形下
部分单位矩阵并非总能得到。
于是我们把所关心的矩阵形式
推广到其自然的一般化，即对角矩阵，
并证明了当特征值互不相同时，
变换或方阵可以对角化。
但我们也给出过一个方阵的例子，
因为它是幂零的，故不能对角化，
所以对角形式不足以充当矩阵相似的规范形式。

上一小节对那个例子作了深入展开，
得到了幂零矩阵的规范形式，即非零元都位于次对角线上的矩阵。

本节将完成我们的计划：证明对任意线性变换都
存在一个基~\( B \)，
使得矩阵表示 $\rep{t}{B,B}$ 是一个对角矩阵
与一个幂零矩阵的和。
这就是 Jordan 标准形。









\subsectionoptional{映射与矩阵的多项式}\index{polynomial!of map, matrix}
回顾一下，方阵的集合~\( \matspace_{\nbyn{n}} \)
在按元素的加法与标量乘法下是一个向量空间，并且
%the unit matrices\Dash all entries are zero except
%for a single entry, which is one\Dash form a basis, so 
这个空间的维数是 \( n^2 \)。
于是对任意 \( \nbyn{n} \) 矩阵 $T$，
含 \( n^2+1 \) 个成员的集合 \( \set{I,T,T^2,\dots,T^{n^2} } \) 线性相关，
因此存在标量 \( c_0,\dots,c_{n^2} \)，
不全为零，使得
\begin{equation*}
  c_{n^2}T^{n^2}+\dots+c_1T+c_0I
\end{equation*}
是零矩阵。
因此每个变换都有一种广义的幂零性\Dash 方阵的幂
不可能无限攀升而不出现某种重复。

\begin{definition}
设 \( t \) 是向量空间~\( V \) 上的线性变换。
若 \( f(x)=c_nx^n+\dots+c_1x+c_0 \) 是一个多项式，
则 \( f(t) \) 是~\( V \) 上的变换
\( c_nt^n+\dots+c_1t+c_0(\identity) \)。
同样地，若 \( T \) 是方阵，
则
\( f(T) \) 是矩阵 \( c_nT^n+\dots+c_1T+c_0I \)。
\end{definition}

\noindent 矩阵的多项式表示映射的多项式：~若
\( T=\rep{t}{B,B} \) 则 \( f(T)=\rep{f(t)}{B,B} \)。
这是因为 \( T^j=\rep{t^j}{B,B} \)，
以及 \( cT=\rep{ct}{B,B} \)，且 \( T_1+T_2 =\rep{t_1+t_2}{B,B} \)。

\begin{remark}
矩阵多项式的写法，我们将与映射多项式的写法略有不同。
例如，若 \( f(x)=x-3 \)，则
我们会写出单位矩阵，如~\( f(T)=T-3I \)，
但不写出恒等映射，如~\( f(t)=t-3 \)。
\end{remark}

\begin{example}  \label{ex:PolySendRotMatToZ}
将平面向量逆时针旋转 \( \pi/3 \)~弧度，
用标准基表示为
\begin{equation*}
  T=
  \begin{mat}[r]
      \cos(\pi/3)  &-\sin(\pi/3)  \\
      \sin(\pi/3)   &\cos(\pi/3)
   \end{mat}
   =
   \begin{mat}[r]
      1/2        &-\sqrt{3}/2  \\
      \sqrt{3}/2  &1/2
   \end{mat}
\end{equation*}
验证 \( T^2-T+I=Z \) 是例行的。
（从几何上看，\( T^2 \) 旋转了 \( 2\pi/3\)。
在标准基向量~\( \vec{e}_1 \) 上，
$T^2-T$ 的作用是给出角为
\( 2\pi/3\) 的单位向量与角为 \( \pi/3\) 的单位向量之差，
即 \( \binom{-1}{0} \)。
在 \( \vec{e}_2 \) 上它给出 \( \binom{0}{-1} \)。
所以 \( T^2-T=-I \)。）
\end{example}

空间 $\matspace_{\nbyn{2}}$ 的维数是四，所以我们知道对任意
\( \nbyn{2} \) 矩阵~\( T \)，存在四次多项式
\( f  \) 使得 \( f(T)=Z \)。
但在前面的例题中我们展示了一个二次的多项式就够用了。
因此虽然次数~$n^2$ 总是够用，但在某些情形
次数更低的多项式就足够了。

\begin{definition}
变换 \( t \)\index{transformation!minimal polynomial}
或方阵 \( T \)\index{matrix!minimal polynomial} 的
\definend{极小多项式}\index{polynomial!minimal}%
\index{minimal polynomial}
\( m(x) \) 是次数最低且首项系数为~\( 1 \)
的非零多项式，
它使得 \( m(t) \) 是零映射或 \( m(T) \) 是零矩阵。
\end{definition}

\noindent 极小多项式不能是常值多项式，因为
对首项系数有限制。
所以极小多项式的
次数至少为一。
零矩阵的极小多项式是 $p(x)=x$，而
单位矩阵的极小多项式是 $\hat{p}(x)=x-1$。

\begin{lemma}
任何变换或方阵都有唯一的极小多项式。  
\end{lemma}

\begin{proof}
先证存在性。
由前面的观察，
次数~$n^2$ 就够用，故至少有一个
非零多项式 $p(x)=c_kx^k+\cdots+c_0$
将该映射或矩阵变为零。
在所有这样的多项式中
取一个次数最小的。
再用其首项系数~$c_k$ 除这个多项式，得到首项为~$1$ 的多项式。
因此任何映射或矩阵至少有一个极小多项式。

下面证唯一性。
假设
\( m(x) \) 与 \( \hat{m}(x) \) 都将该映射或矩阵变为零，
二者都是
次数最小的，从而次数相等，
且首项都是~$1$。
考虑差 \( m(x)-\hat{m}(x) \)。
若它不是零多项式，则它有非零的首项系数。
用这个首项系数
去除整个多项式，将得到一个多项式，它
把该映射或矩阵变为零，有首项
系数~\( 1\)，且次数比
$m$ 与 $\hat{m}$ 都低
（因为相减时首项的 \( 1\) 相消）。
这会与 $m$ 和 $\hat{m}$ 的次数
的最小性矛盾。
于是 \( m(x)-\hat{m}(x) \) 是零多项式，
即二者相等。
\end{proof}

\begin{example}  \label{ex:MinPolyForRotMat}
计算\nearbyexample{ex:PolySendRotMatToZ} 中矩阵的
极小多项式的一种方法，是求出 $T$ 的幂
直到 $n^2=4$。
\begin{equation*}
    T^2=
    \begin{mat}[r]
       -1/2         &-\sqrt{3}/2  \\
       \sqrt{3}/2   &-1/2
    \end{mat} 
    \quad
    T^3=
    \begin{mat}[r]
       -1    &0           \\
        0    &-1
    \end{mat}
    \quad
    T^4=
    \begin{mat}[r]
       -1/2         &\sqrt{3}/2  \\
       -\sqrt{3}/2  &-1/2
    \end{mat}
\end{equation*}
令 \( c_4T^4+c_3T^3+c_2T^2+c_1T+c_0I \) 等于零矩阵。
\begin{equation*}
  \begin{linsys}{5}
      -(1/2)c_4  
          &-  &c_3           
          &- &(1/2)c_2
          &+ &(1/2)c_1  
          &+  &c_0  &=  &0      \\
      (\sqrt{3}/2)c_4  
          &  &   
          &- &(\sqrt{3}/2)c_2
          &- &(\sqrt{3}/2)c_1  
          &   &     &=  &0        \\
      -(\sqrt{3}/2)c_4  
          &   &    
          &+ &(\sqrt{3}/2)c_2
          &+ &(\sqrt{3}/2)c_1  
         &   &            &=  &0      \\
      -(1/2)c_4  
          &-  &c_3             
          &- &(1/2)c_2
          &+ &(1/2)c_1  
          &+  &c_0  &=  &0
    \end{linsys}
\end{equation*}
用高斯消元法。
\begin{equation*}
  \begin{linsys}{5}
      -(1/2)c_4  
          &- &c_3             
          &- &(1/2)c_2
          &+ &(1/2)c_1  
          &+  &c_0  &=  &0      \\
          &  &-(\sqrt{3})c_3 
          &- &\sqrt{3}c_2
          &  &    
          &+ &\sqrt{3}c_0 &=  &0
    \end{linsys} 
\end{equation*}
着眼于使多项式的次数尽可能
小，注意
令 \( c_4 \)、\( c_3 \) 与 \( c_2 \) 为零会迫使 \( c_1 \) 与
\( c_0 \) 也得出为零，
所以这些方程不允许一次的极小多项式。
于是改为令 \( c_4 \) 与 \( c_3 \) 为零。
方程组
\begin{equation*}
  \begin{linsys}{3}           
          -(1/2)c_2
          &+ &(1/2)c_1  
          &+  &c_0  &=  &0      \\
          -\sqrt{3}c_2
          &  &    
          &+ &\sqrt{3}c_0 &=  &0
    \end{linsys} 
\end{equation*}
的解集为 \( c_1=-c_0 \) 与 \( c_2=c_0 \)。
取首项系数为~\( c_2=1 \)
得到极小多项式 $x^2-x+1$。
\end{example}

% Using the method of that example to find the minimal polynomial of a
% \( \nbyn{3} \) matrix 
% would mean doing Gaussian reduction on
% a system with nine equations in ten unknowns.
那个计算很笨拙。
我们将给出另一种方法。
% To begin, note that we can break a polynomial of a map or a matrix into
% its components.
% (For this lemma, recall that we are using complex numbers in this chapter
% so all polynomials break completely into linear factors.)

\begin{lemma} \label{le:PolyMapsFactor}
假设多项式 \( f(x)=c_nx^n+\dots+c_1x+c_0 \) 可以因式分解为
\( k(x-\lambda_1)^{q_1}\cdots(x-\lambda_z)^{q_z} \)。
若 \( t \) 是线性变换，则下面两个是相等的映射。
\begin{equation*}
  c_nt^n+\dots+c_1t+c_0
  =
  k\cdot\composed{\composed{(t-\lambda_1)^{q_1}}{\cdots}}{
      (t-\lambda_z)^{q_z}} 
\end{equation*}
因此，若 \( T \) 是方阵，则 \( f(T) \) 与
\( k\cdot(T-\lambda_1I)^{q_1}\cdots(T-\lambda_z I)^{q_z} \) 
是相等的矩阵。
\end{lemma}

\begin{proof}
我们对多项式的次数用归纳法。
多项式次数为~零与次数为~一的情形是显然的。
完整的归纳论证是\nearbyexercise{exer:PolyMapsFactor}，
但我们将以次数~二的情形给出其大意。

二次多项式可以分解成两个
一次项 \( f(x)=k(x-\lambda_1)\cdot(x-\lambda_2)
                    =k(x^2+(-\lambda_1-\lambda_2)x+\lambda_1\lambda_2) \)。
% (the roots $\lambda_1$ and $\lambda_2$ could be equal).
代入 \( t \) 
——无论是分解后的还是未分解的形式——得到的都是同一个映射。
\begin{align*}
   \bigl(k\cdot\composed{(t-\lambda_1)}{(t-\lambda_2)}\bigr)\>(\vec{v})
   &=\bigl(k\cdot(t-\lambda_1)\bigr)\,(t(\vec{v})-\lambda_2\vec{v})    \\
   &=k\cdot\bigl(t(t(\vec{v}))-t(\lambda_2\vec{v})
      -\lambda_1 t(\vec{v})+\lambda_1\lambda_2\vec{v}\bigr)    \\
   &=k\cdot \bigl(\composed{t}{t}\,(\vec{v})-(\lambda_1+\lambda_2)t(\vec{v})
          +\lambda_1\lambda_2\vec{v}\bigr)                    \\
   &=k\cdot(t^2-(\lambda_1+\lambda_2)t+\lambda_1\lambda_2)\>(\vec{v})
\end{align*}
第三个等号利用 \( t \) 的线性性把 $\lambda_2$ 从
第二项中提出来。
\end{proof}

% In particular, if a minimal polynomial for a transformation $t$ 
% factors as
% $m(x)=(x-\lambda_1)^{q_1}\cdots (x-\lambda_z)^{q_z}$
% then 
% \( m(t)=\composed{\composed{(t-\lambda_1)^{q_1}}{\cdots}}{
%       (t-\lambda_z)^{q_z}} \) 
% is the zero map. 
% Since \( m(t) \) sends every vector to zero, at least
% one of the maps \( t-\lambda_i \)  sends some
% nonzero vectors to zero.
% Exactly the same holds in the matrix case\Dash if $m$ is minimal for $T$ then
% \( m(T)=(T-\lambda_1I)^{q_1}\cdots (T-\lambda_z I)^{q_z} \)
% is the zero matrix and at least one of the matrices $T-\lambda_iI$
% sends some nonzero vectors to zero. 
% That is, in both cases at least some of the \( \lambda_i \) are eigenvalues.
% (\nearbyexercise{exer:SomeRootsMinPolyAreEigs} expands on this.)

下一个结论是：
极小多项式的每个根都是一个特征值，
且每个特征值都是极小多项式的根。
（也就是说，该结论说的是 `$1\leq q_i$' 而不只是
`$0\leq q_i$'。）
为此，回顾求特征值时
我们求解 $\deter{T-xI}=0$，
这个行列式给出一个关于 $x$ 的多项式，
即特征多项式，
它的根就是特征值。

\begin{theorem}[Cayley-Hamilton]
\label{th:CayHam}
\index{Cayley-Hamilton theorem}
\hspace*{0em plus2em}
若变换或方阵的特征多项式
分解为
\begin{equation*}
  k\cdot (x-\lambda_1)^{p_1}(x-\lambda_2)^{p_2}\cdots(x-\lambda_z)^{p_z}
\end{equation*}
则其极小多项式分解为
\begin{equation*}
  (x-\lambda_1)^{q_1}(x-\lambda_2)^{q_2}\cdots(x-\lambda_z)^{q_z}
\end{equation*}
其中对 \( 1 \) 到 \( z \) 之间的每个 \( i \) 有 \( 1\leq q_i \leq p_i \)。
\end{theorem}

\noindent 证明占接下来的三个引理。
我们将用矩阵的术语来叙述它们，
因为这个版本更方便
第一个证明，但它们同样
适用于映射。

第一个结论是关键。
为证明它，注意我们可以把
一个元素为多项式的矩阵看成
系数为矩阵的多项式。
\begin{equation*}
   \begin{mat}
     2x^2+3x-1  &x^2+2    \\
     3x^2+4x+1  &4x^2+x+1
   \end{mat}
 = \begin{mat}[r]
    2  &1  \\
    3  &4
  \end{mat}x^2
 + \begin{mat}[r]
    3  &0  \\
    4  &1
  \end{mat}x
 + \begin{mat}[r]
   -1  &2  \\
    1  &1
  \end{mat}
\end{equation*}

\begin{lemma}   \label{le:MatSatItsCharPoly}
若 \( T \) 是特征多项式为 \( c(x) \) 的方阵，
则 \( c(T) \) 是零矩阵。
\end{lemma}

\begin{proof}
设 \( C \) 为 \( T-xI \)，
它的行列式即特征多项式
\( c(x)=c_nx^n+\dots+c_1x+c_0 \)。
\begin{equation*}
  C=\begin{mat}
    t_{1,1}-x        &t_{1,2}   &\ldots        \\
    t_{2,1}          &t_{2,2}-x               \\
    \vdots           &          &\ddots       \\
                     &          &       &t_{n,n}-x
  \end{mat}
\end{equation*}
回顾定理~Four.III.\ref{th:MatTimesAdjEqDiagDets}：
矩阵与其伴随矩阵的乘积等于
矩阵的行列式乘以单位矩阵。
\begin{equation*}
  c(x)\cdot I
  =\adj (C)C
  =\adj (C)(T-xI)
  =\adj (C)T- \adj(C)\cdot x
\tag*{($*$)}
\end{equation*}
~($*$) 的左边是
$c_nIx^n+c_{n-1}Ix^{n-1}+\cdots+c_1Ix+c_0I$。
对于右边，
\( \adj (C) \) 的元素是多项式，每个的次数
至多是 \( n-1 \)，因为矩阵的子式去掉了一行与一列。
正如证明之前所提示的，把它改写成一个
系数为矩阵的多项式：
\( \adj (C)=C_{n-1}x^{n-1}+\cdots+C_1x+C_0 \)，
其中每个 \( C_i \) 是标量组成的矩阵。
现在~($*$) 的右边是
\begin{equation*}
  [(C_{n-1}T)x^{n-1}+\cdots+(C_1T)x+C_0T]  
   -[C_{n-1}x^n+C_{n-2}x^{n-1}+\cdots+C_0x]
\end{equation*}
比较 ($*$) 左右两边
\( x^n \)、$x^{n-1}$ 等的
系数。
\begin{align*}
  c_nI
  &=-C_{n-1}    \\
  c_{n-1}I
  &=-C_{n-2}+C_{n-1}T    \\
  &\vdotswithin{=}             \\
  c_{1}I
  &=-C_{0}+C_{1}T    \\
  c_{0}I
  &=C_{0}T
\end{align*}
将第一个等式的两边从右边乘以 \( T^n \)，
第二个等式的两边乘以 \( T^{n-1} \)，等等。
\begin{align*}
  c_nT^n
  &=-C_{n-1}T^n    \\
  c_{n-1}T^{n-1}
  &=-C_{n-2}T^{n-1}+C_{n-1}T^n    \\
  &\vdotswithin{=}             \\
  c_{1}T
  &=-C_{0}T+C_{1}T^2    \\
  c_{0}I
  &=C_{0}T
\end{align*}
相加。
左边是
\( c_nT^n+c_{n-1}T^{n-1}+\cdots+c_0I \)。
右边是叠缩的；
例如，第一行的 $-C_{n-1}T^n$ 与第二行的
一半 $C_{n-1}T^n$ 相消。
右边总共是零矩阵。
\end{proof}

我们说矩阵或映射
\definend{满足}\index{characteristic!polynomial!satisfied by} 
其特征多项式，指的就是那个结论。

\begin{lemma} \label{le:tSatisImpMinPolyDivides}
任何被 \( T \) 满足的多项式都可被
\( T \) 的极小多项式整除。
也就是说，对多项式 \( f(x) \)，若 \( f(T) \) 是零矩阵，
则 \( f(x) \) 可被 \( T \) 的极小多项式整除。
\end{lemma}

\begin{proof}
设 \( m(x) \) 是 \( T \) 的极小多项式。
多项式的带余除法定理给出
\( f(x)=q(x)\cdot m(x)+r(x) \)，
其中 \( r \) 的次数严格小于 \( m \) 的次数。
因为 $T$ 满足 $f$ 与 $m$，把 $T$ 代入那个等式得到
\( r(T) \) 是零矩阵。
除非 \( r \)
是零多项式，否则这与 \( m \) 的极小性矛盾。
\end{proof}

把前面两个引理合起来，可知极小多项式
整除特征多项式。
于是
极小多项式的任意根也是特征
多项式的根。

于是到目前为止我们有：若极小多项式分解为
\( m(x)=(x-\lambda_1)^{q_1}\cdots(x-\lambda_i)^{q_i} \)，则
特征多项式也有根 \( \lambda_1\), \ldots, 
\( \lambda_i \)。
但就我们至今所建立的而言，
特征多项式可能还有别的根：
\( c(x)=(x-\lambda_1)^{p_1}\cdots(x-\lambda_i)^{p_i}
     (x-\lambda_{i+1})^{p_{i+1}}\cdots(x-\lambda_z)^{p_z} \)，其中
对 \( 1\leq j \leq i \) 有 \( 1\leq q_j\leq p_j \)。
我们通过证明特征多项式没有另外的根来
完成 Cayley-Hamilton 定理的证明，于是
不存在 $\lambda_{i+1}$、$\lambda_{i+2}$ 等。

\begin{lemma}
方阵的特征多项式的每个一次因式
也是极小多项式的一次因式。
\end{lemma}

\begin{proof}
设 \( T \) 是极小多项式 \( m(x) \)
次数为~\( n \) 的方阵，
并假设 \( x-\lambda \) 是 \( T \) 的特征多项式的一个因式，
从而 \( \lambda \) 是 \( T \) 的一个特征值。
我们将证明 $m(\lambda)=0$，即 $x-\lambda$ 是 $m$ 的因式。

假设 $\lambda$ 与特征向量~$\vec{v}$ 相联系，
考虑幂 $T^2(\vec{v})$, \ldots, \( T^n(\vec{v}) \)。
我们有
$T^2(\vec{v})
=T\cdot\lambda\vec{v}=\lambda\cdot T\vec{v}=\lambda^2\vec{v}$。
所有的幂都是如此：
对 \( 1\leq i \leq n \) 有 $T^i\vec{v}=\lambda^i\vec{v}\,$。
于是
对任意多项式函数 \( p(x) \)，把矩阵 \( p(T) \)
作用到 \( \vec{v} \) 上等于用标量
\( p(\lambda) \) 乘 \( \vec{v} \) 的结果。
\begin{multline*}
  p(T)\,\vec{v}
  =(c_kT^k+\dots+c_1T+c_0I)\,\vec{v}
  =c_kT^k\vec{v}+\dots+c_1T\vec{v}+c_0\vec{v}  \\
  =c_k\lambda^k\vec{v}+\dots+c_1\lambda\vec{v}+c_0\vec{v}
  =p(\lambda)\cdot\vec{v}
\end{multline*}
由于 \( m(T) \) 是零矩阵，
故 \( \zero=m(T)(\vec{v})=m(\lambda)\cdot\vec{v} \)，
从而 \( m(\lambda)=0 \)，因为 \( \vec{v}\neq\zero \)
（它是特征向量）。
\end{proof}

Cayley-Hamilton 定理的证明到此结束。

\begin{example} \label{ex:MinPolyUsingCH}
我们可以用 Cayley-Hamilton 定理来求
这个矩阵的极小多项式。
\begin{equation*}
   T=
   \begin{mat}[r]
      2  &0  &0  &1  \\
      1  &2  &0  &2  \\
      0  &0  &2  &-1 \\
      0  &0  &0  &1
   \end{mat}
\end{equation*}
先用通常的行列式 $\deter{T-xI}$ 求出它的特征多项式 \( c(x)=(x-1)(x-2)^3 \)。
有了它，Cayley-Hamilton 定理就说
\( T \) 的极小多项式是
\( (x-1)(x-2) \)、
\( (x-1)(x-2)^2 \) 或
\( (x-1)(x-2)^3 \) 之一。
只需计算
\begin{equation*}
   (T-1I)(T-2I)=\!
   \begin{mat}[r]
      1  &0  &0  &1  \\
      1  &1  &0  &2  \\
      0  &0  &1  &-1 \\
      0  &0  &0  &0
   \end{mat}
   \begin{mat}[r]
      0  &0  &0  &1  \\
      1  &0  &0  &2  \\
      0  &0  &0  &-1 \\
      0  &0  &0  &-1
   \end{mat}
   =
   \begin{mat}[r]
      0  &0  &0  &0  \\
      1  &0  &0  &1  \\
      0  &0  &0  &0  \\
      0  &0  &0  &0
   \end{mat}
\end{equation*}
与
\begin{equation*}
   (T-1I)(T-2I)^2=
   \begin{mat}[r]
      0  &0  &0  &0  \\
      1  &0  &0  &1  \\
      0  &0  &0  &0  \\
      0  &0  &0  &0
   \end{mat}
   \begin{mat}[r]
      0  &0  &0  &1  \\
      1  &0  &0  &2  \\
      0  &0  &0  &-1 \\
      0  &0  &0  &-1
   \end{mat}
   =
   \begin{mat}[r]
      0  &0  &0  &0  \\
      0  &0  &0  &0  \\
      0  &0  &0  &0  \\
      0  &0  &0  &0
   \end{mat}
\end{equation*}
于是 \( m(x)=(x-1)(x-2)^2 \)。
\end{example}


% ==================================
